% yquant circuit diagrams; \input by main.tex after \usepackage{yquant}.
\usetikzlibrary{calc,decorations.pathreplacing}
\yquantset{
  register/minimum height=3.6mm,
  register/separation=1.8mm,
  operator/separation=0.9mm,
  operator/minimum width=4mm,
  operators/every barrier/.append style={thin, dashed, gray},
}
\tikzset{
  stage/.style={decorate, decoration={brace, amplitude=3pt, raise=1pt},
    gray!70!black},
  stage label/.style={font=\footnotesize, gray!70!black, above=3pt, midway},
}
% Brace above the circuit from the west of node #1 to the east of node #2.
\newcommand{\circstage}[3]{%
  \draw[stage] (#1.west |- top) -- (#2.east |- top)
    node[stage label] {#3};}
\newcommand{\circphase}[1]{$\scriptstyle #1$}

% Adapted-basis circuit in the field encoding, drawn for n=3 (p=5,7).
\newcommand{\AdaptedCircuit}{%
\begin{tikzpicture}
\begin{yquant}
qubit {$x_{\idx}$} x[3];
qubit {$a_{\idx}$} a[3];
qubit {$y_{\idx}$} w[3];

[name=prep] box {$U_p$} (x);
box {$V_q$} (a);
box {$U_p$} (w);
barrier (-);
[name=o1] box {\circphase{\theta_0}} w[0] | x[0];
box {\circphase{\theta_1}} w[1] | x[0];
box {\circphase{\theta_2}} w[2] | x[0];
box {\circphase{\theta_1}} w[0] | x[1];
box {\circphase{\theta_2}} w[1] | x[1];
box {\circphase{\theta_3}} w[2] | x[1];
box {\circphase{\theta_2}} w[0] | x[2];
box {\circphase{\theta_3}} w[1] | x[2];
[name=o2] box {\circphase{\theta_4}} w[2] | x[2];
[name=o3] box {\circphase{b\theta_0}} w[0] | a[0];
box {\circphase{b\theta_1}} w[1] | a[0];
box {\circphase{b\theta_2}} w[2] | a[0];
box {\circphase{b\theta_1}} w[0] | a[1];
box {\circphase{b\theta_2}} w[1] | a[1];
box {\circphase{b\theta_3}} w[2] | a[1];
box {\circphase{b\theta_2}} w[0] | a[2];
box {\circphase{b\theta_3}} w[1] | a[2];
[name=o4] box {\circphase{b\theta_4}} w[2] | a[2];
barrier (-);
[name=ft] box {$QFT_p$} (x);
discard w;
measure x;
[name=fa] box {$QFT_p$} (a) | x;
[name=m] measure a;
output {$k$} (x);
output {$l$} (a);
\end{yquant}
\node[font=\scriptsize, anchor=south west, inner sep=1pt]
  at ($(fa.north)+(0.6mm,0.4mm)$) {$k\neq0$};
\coordinate (top) at ($(current bounding box.north)+(0,1mm)$);
\circstage{prep}{prep}{prepare}
\circstage{o1}{o2}{oracle, $x$ on $y$}
\circstage{o3}{o4}{oracle, $a$ on $y$}
\circstage{ft}{m-0}{QFT, measure}
\end{tikzpicture}}

% Exponent-encoding circuit with registers drawn as bundles.
\newcommand{\ExponentCircuit}{%
\begin{tikzpicture}
\begin{yquant}
qubit {$x$} x;
qubit {$t$} t;
qubit {$y$} w;
slash x; slash t; slash w;
hspace {2mm} x; hspace {2mm} t; hspace {2mm} w;
box {$U_p$} x;
box {$U_q$} t;
box {$U_p$} w;
box {$\Theta$} w | x;
box {$D_b$} (t, w);
box {$QFT_p$} x;
box {$M$} t;
discard w;
measure x;
measure t;
output {$k$} x;
output {$s$} t;
\end{yquant}
\end{tikzpicture}}

% One measurement variant M on t = (t_0, t_{>0}); #1 is the yquant body.
\newcommand{\Variant}[2]{%
\begin{tabular}[b]{@{}c@{}}
\begin{tikzpicture}[baseline=(current bounding box.center)]
\begin{yquant}
qubit {$t_0$} t0;
qubit {$t_{>0}$} t1;
slash t1;
hspace {2mm} t0; hspace {2mm} t1;
#1
hspace {1.5mm} -;
\end{yquant}
\end{tikzpicture}\\[1pt]
\footnotesize #2
\end{tabular}}

\newcommand{\MeasurementVariants}{%
\Variant{box {$T$} t0; box {$H$} t0;}{pairs}\hfill
\Variant{box {$H$} t0;}{$H$ only}\hfill
\Variant{box {$F_q$} (t0, t1);}{forgetful}\hfill
\Variant{hspace {5mm} -;}{weak}}

% ZNE folding of one native ZZ gate, as in forte.fold_zz.
\newcommand{\FoldCircuit}{%
\begin{tikzpicture}[baseline=(current bounding box.center)]
\begin{yquant}
qubit {} q[2];
hspace {1mm} q;
box {$Z\!Z(\theta)$} (q);
hspace {1mm} q;
\end{yquant}
\end{tikzpicture}
\;$\longrightarrow$\;
\begin{tikzpicture}[baseline=(current bounding box.center)]
\begin{yquant}
qubit {} q[2];
hspace {1mm} q;
box {$Z\!Z(\theta)$} (q);
box {$X$} q[0];
box {$Z\!Z(\theta)$} (q);
box {$X$} q[0];
box {$Z\!Z(\theta)$} (q);
hspace {1mm} q;
\end{yquant}
\end{tikzpicture}}


% q=6: two strings per w_j cancelled; one CNOT walk through t_0 serves every w_j.
\newcommand{\OracleDiagonalSix}{%
\begin{tikzpicture}[baseline=(current bounding box.center), scale=0.8]
\begin{yquant}
qubit {$t_{\idx}$} t[3];
qubit {$y_{\idx}$} w[3];
[name=s1a] zz (t[0], w[0]);
zz (t[0], w[1]);
[name=s1b] zz (t[0], w[2]);
box {$\scriptstyle R_Z$} t[0];
[name=s2a] zz (t[1], w[0]);
zz (t[1], w[1]);
[name=s2b] zz (t[1], w[2]);
[name=s3a] zz (t[2], w[0]);
zz (t[2], w[1]);
[name=s3b] zz (t[2], w[2]);
cnot t[1] | t[0];
[name=s4a] zz (t[1], w[0]);
zz (t[1], w[1]);
[name=s4b] zz (t[1], w[2]);
cnot t[1] | t[0];
cnot t[2] | t[0];
[name=s5a] zz (t[2], w[0]);
zz (t[2], w[1]);
[name=s5b] zz (t[2], w[2]);
cnot t[2] | t[0];
\end{yquant}
\coordinate (top) at ($(current bounding box.north)+(0,1mm)$);
\draw[stage] (s1a.west |- top) -- (s1b.east |- top) node[stage label] {$t_0$};
\draw[stage] (s2a.west |- top) -- (s2b.east |- top) node[stage label] {$t_1$};
\draw[stage] (s3a.west |- top) -- (s3b.east |- top) node[stage label] {$t_2$};
\draw[stage] (s4a.west |- top) -- (s4b.east |- top) node[stage label] {$t_0{\oplus}t_1$};
\draw[stage] (s5a.west |- top) -- (s5b.east |- top) node[stage label] {$t_0{\oplus}t_2$};
\end{tikzpicture}}
